\section*{Chapitre \ref{ch:g9:landforms} --- Les formes du relief}

\begin{solution}{exo:g9:landforms:1}
(a) Le vent. (b) L'eau qui dissout la \omterm{def:g3:stones-pebbles-sand:rock}{roche} (et qui dépose de la calcite
dans la \omterm{def:g9:landforms:karst}{grotte}). (c) La \omterm{def:g2:water-on-earth:ocean}{mer} (\omterm{def:g9:landforms:beach}{dérive littorale}). (d) La glace. (e) La \omterm{def:g2:water-on-earth:ocean}{mer}.
(f) Une \omterm{def:g1:landscapes-around-us:water}{rivière}.
\end{solution}

\begin{solution}{exo:g9:landforms:2}
(a) Le \omterm{def:g9:landforms:uplift}{soulèvement}. (b) La \omterm{def:g9:landforms:dune}{face d'avalanche}. (c) Une \omterm{def:g9:landforms:karst}{doline}. (d) La \omterm{def:g9:landforms:beach}{dérive
littorale}. (e) Un \omterm{def:g9:landforms:karst}{karst}.
\end{solution}

\begin{solution}{exo:g9:landforms:3}
(a) Faux~: les cornes pointent dans le sens du vent~; le nez arrondi fait
face au vent. (b) Vrai. (c) Faux~: le vent la déplace grain par grain, en
remontant la pente douce et en descendant la \omterm{def:g9:landforms:dune}{face d'avalanche}. (d) Vrai.
\end{solution}

\begin{solution}{exo:g9:landforms:4}
(b), (d), (e), (a), (c).
\end{solution}

\begin{solution}{exo:g9:landforms:5}
(a) De l'ouest~: les cornes pointent dans le sens du vent. (b) Du côté
est, entre les cornes. (c) $60 \div 6 = 10$ ans. (d) Une \omterm{def:g9:landforms:dune}{dune} plus grande
contient plus de \omterm{def:g3:stones-pebbles-sand:sand}{sable}, et tout ce \omterm{def:g3:stones-pebbles-sand:sand}{sable} doit passer par-dessus la crête
pour que la \omterm{def:g9:landforms:dune}{dune} avance.
\end{solution}

\begin{solution}{exo:g9:landforms:6}
(a) $7 \times 1 = \qty{7}{m}$. (b) $2 \div 12 = \frac{1}{6}$ d'année.
(c) Sept ans de recul en $\frac{1}{6}$ d'année~: $7 \times 6 = 42$
fois plus vite. (d) Les vagues creusent d'abord une encoche~; la \omterm{def:g3:stones-pebbles-sand:rock}{roche}
au-dessus la surplombe pendant un certain temps, sans bouger~; puis le
surplomb s'effondre d'un coup, et le bord recule brusquement.
\end{solution}

\begin{solution}{exo:g9:landforms:7}
(a) Vers l'est. (b) Trois montées sur la \omterm{def:g9:landforms:beach}{plage}, obliques vers le
nord-est, suivies chacune d'un retour tout droit vers le sud~: un zigzag
qui avance vers l'est. (c) Une \omterm{def:g9:landforms:beach}{flèche littorale}, qui s'allonge vers l'est
en travers de l'entrée de la baie. (d) La flèche abrite des vagues l'eau
qui se trouve derrière elle~; dans une eau calme, la vase fine peut se
déposer.
\end{solution}

\begin{solution}{exo:g9:landforms:8}
(a) Elle se charge de dioxyde de carbone dans l'air et dans le sol. (b) Le
faible acide dissout lentement le calcaire, mais n'attaque presque pas le
granite. (c) L'eau ne peut pénétrer dans la \omterm{def:g3:stones-pebbles-sand:rock}{roche} que par ses fissures et
le long de ses couches~; ce sont ces ouvertures qu'elle élargit d'abord.
(d) La \omterm{def:g3:stones-pebbles-sand:rock}{roche} sous la surface est dissoute et le sol s'affaisse lentement~;
ou bien la voûte d'une \omterm{def:g9:landforms:karst}{grotte} s'effondre.
\end{solution}

\begin{solution}{exo:g9:landforms:9}
(a) Quand la goutte rencontre l'air de la \omterm{def:g9:landforms:karst}{grotte}, elle libère une partie
de son dioxyde de carbone et ne peut plus garder tout le calcaire dissous,
qui se dépose sous forme de calcite. (b) Une stalagmite grandit à partir
de gouttes qui s'écrasent sur le sol et s'étalent~; aucune goutte n'est
suspendue à sa pointe pour construire un tube.
(c) $\qty{30}{cm} = \qty{300}{mm}$~; $300 \div 0.1 = 3000 = 3 \times 10^3$
ans.
\end{solution}

\begin{solution}{exo:g9:landforms:10}
(a) $4600 \div 10^7 = 4.6 \times 10^{-4}$~m par an, soit \qty{0.46}{mm}
par an. (b) $0.46 - 0.2 = \qty{0.26}{mm}$ par an. (c) $\qty{1}{m} =
\qty{1000}{mm}$~; $1000 \div 0.46 \approx 2200$, environ deux mille fois
plus vite. (d) Il est maintenu pendant des millions d'années~: $\qty{0.26}{mm} \times
10^7 = 2.6 \times 10^6$~mm, environ \qty{2.6}{km}.
\end{solution}

\begin{solution}{exo:g9:landforms:11}
(a) Forme~: des \omterm{def:g9:landforms:dune}{dunes} en croissant dont la \omterm{def:g9:landforms:dune}{face d'avalanche} est du côté
sud-est. Agent et processus~: le vent, qui souffle du nord-ouest~; \omterm{def:g5:water-shapes-the-land:deposition}{dépôt}.
Matériau~: \omterm{def:g3:stones-pebbles-sand:sand}{sable} meuble. (b) Forme~: une \omterm{def:g9:landforms:cliff}{falaise} au-dessus d'un replat
découvert à marée basse. Agent et processus~: la \omterm{def:g2:water-on-earth:ocean}{mer}, \omterm{def:g5:water-shapes-the-land:erosion}{érosion}~; le replat
est une plate-forme d'abrasion laissée par le recul de la \omterm{def:g9:landforms:cliff}{falaise}.
Matériau~: une \omterm{def:g3:stones-pebbles-sand:rock}{roche} blanche, probablement de la craie, une \omterm{def:g3:stones-pebbles-sand:rock}{roche} tendre.
(c) Forme~: un ruisseau qui disparaît dans un trou, une \omterm{def:g6:water-underground:spring}{source} au loin.
Agent et processus~: l'eau qui dissout la \omterm{def:g3:stones-pebbles-sand:rock}{roche} et circule sous terre.
Matériau~: du calcaire~; le plateau est un \omterm{def:g9:landforms:karst}{karst}.
\end{solution}

\begin{solution}{exo:g9:landforms:12}
(a) La petite \omterm{def:g9:landforms:dune}{dune}~: $45 \div 6 = 7.5$ ans~; la plus grande~: $45 \div 3 =
15$ ans. (b) $45 \div 100 = 0.45$ d'année, environ cinq ou six mois.
(c) Par exemple, planter des herbes sur les \omterm{def:g9:landforms:dune}{dunes}~: leurs racines
retiennent le \omterm{def:g3:stones-pebbles-sand:sand}{sable}, si bien que le vent ne peut plus l'emporter.
\end{solution}

\begin{solution}{pb:g9:landforms:1}
\textbf{1.} $46.0 - 30.4 = \qty{15.6}{m}$.

\textbf{2.} $15.6 \div 20 = \qty{0.78}{m}$ par an.

\textbf{3.} Entre 2012 et 2014~: $41.0 - 36.2 = \qty{4.8}{m}$ en deux
ans, soit $4.8 \div 2 = \qty{2.4}{m}$ par an.

\textbf{4.} $2.4 \div 0.78 \approx 3$ fois plus vite. Probablement un
hiver de \omterm{def:g5:natural-hazards:storm}{tempêtes}, pendant lequel de grands surplombs se sont effondrés.

\textbf{5.} \qty{78}{cm} par an~; $0.78 \times 100 = \qty{78}{m}$ par
siècle.

\textbf{6.} Oui~: \qty{0.78}{m} par an, c'est un peu moins d'un mètre par an.

\textbf{7.} $17 \div 0.78 \approx 22$ ans.

\textbf{8.} La \omterm{def:g9:landforms:cliff}{falaise} recule par à-coups~: un intervalle de deux ans peut
contenir un grand éboulement, ou aucun. Sur 20 ans, les à-coups et les
années calmes se compensent.

\textbf{9.} Le bord a encore $30.4 - 10 = \qty{20.4}{m}$ à parcourir~: $20.4 \div 0.78
\approx 26$ ans, donc en 2050.

\textbf{10.} $2 \times 0.78 = \qty{1.56}{m}$ par an~; $30.4 \div 1.56
\approx 19.5$ ans après 2024, donc en 2043.

\textbf{11.} Le bord d'une \omterm{def:g9:landforms:cliff}{falaise} peut céder à tout moment~; se tenir
dessus, ou tout près, est dangereux.

\textbf{12.} $30.4 \div 0.78 \approx 39$ ans après 2024~: en
\textbf{2063}.
\end{solution}
